% Generated draft: requires mathematical and source review.
\documentclass[11pt]{article}
\usepackage[margin=25mm]{geometry}
\usepackage{fontspec}
\setmainfont{Latin Modern Roman}
\usepackage{amsmath,amssymb,mathtools,unicode-math}
\setmathfont{Latin Modern Math}
\usepackage{enumitem,needspace,xurl,hyperref}
\setlength{\parindent}{0pt}
\setlength{\parskip}{5pt}
\setlength{\emergencystretch}{3em}
\newcommand{\N}{\mathbb{N}}
\newcommand{\Z}{\mathbb{Z}}
\newcommand{\Q}{\mathbb{Q}}
\newcommand{\R}{\mathbb{R}}
\newcommand{\C}{\mathbb{C}}
\newcommand{\F}{\mathbb{F}}
\newcommand{\A}{\mathbb{A}}
\newcommand{\PP}{\mathbb{P}}
\newcommand{\E}{\mathbb{E}}
\newcommand{\eps}{\varepsilon}
\newcommand{\dd}{\,\mathrm{d}}
\begin{document}
\section*{UnsolvedMath ID: 30001189}
\textbf{Source problem code:} catalogue entry
\subsection{Definitions and mathematical statement}
Let \(\mathcal T\) be a \(K\)-linear triangulated category with suspension functor \([1]\), and let \((\mathcal T^{\leq 0},\mathcal T^{\geq 0})\) be a \(t\)-structure with heart
\[
\mathcal A=\mathcal T^{\leq 0}\cap\mathcal T^{\geq 0}.
\]
The source does not specify whether \(K\) is a field or a more general coefficient ring, or impose boundedness on the \(t\)-structure.

For \(n\in\Z\), define the degree-\(n\) graded centre by
\[
Z^n(\mathcal T)=\left\{\eta:\operatorname{Id}_{\mathcal T}\Rightarrow[n]\;\middle|\;\eta_{Y[1]}=(-1)^n\eta_Y[1]\text{ for every }Y\in\mathcal T\right\}.
\]
Here \(\eta\) is a natural transformation and the displayed equality uses the usual identifications of iterated shifts. Set \(Z^*(\mathcal T)=\bigoplus_{n\in\Z}Z^n(\mathcal T)\), with multiplication given by composition after shifting. Write \(HH^*(\mathcal A)\) for the Hochschild cohomology of the abelian category \(\mathcal A\) intended by the source. Its construction, coefficient conventions, and the hypotheses needed for it to be defined are not specified in the supplied record.

\textbf{Question.} For every such \(\mathcal T\) and heart \(\mathcal A\), is there a characteristic homomorphism of graded algebras
\[
\chi_{\mathcal A,\mathcal T}:HH^*(\mathcal A)\longrightarrow Z^*(\mathcal T),
\]
and is this construction independent of the choice of heart? In particular, a homogeneous class \(a\in HH^n(\mathcal A)\) would give natural morphisms
\[
\chi_{\mathcal A,\mathcal T}(a)_Y:Y\longrightarrow Y[n]
\]
satisfying the graded-centre sign condition. The intended characteristic or canonical property is not specified by the record; mere existence of an arbitrary graded algebra homomorphism does not express this additional requirement. Moreover, for two hearts \(\mathcal A\) and \(\mathcal B\) in the same \(\mathcal T\), the domains \(HH^*(\mathcal A)\) and \(HH^*(\mathcal B)\) differ. The source supplies no comparison maps or identifications under which heart independence could be stated as an equality or a commuting diagram.

\textbf{Geometric question.} For an analytic space or scheme \(X\), the source asks for the analogous construction, replacing \(D(\mathcal A)\) by \(D(X)\), replacing \(\mathcal A\) by the category of quasi-coherent sheaves, and using \(HH^*(X)\). Its intended map has the form
\[
\chi_X:HH^*(X)\longrightarrow Z^*(D(X)),
\]
with the analogous independence question. Here \(D(\mathcal A)\) denotes a derived category of \(\mathcal A\); the source does not fix its boundedness convention. It also does not specify the sheaf category or boundedness defining \(D(X)\), the construction of \(HH^*(X)\), the hypotheses on \(X\), or the analytic meaning of the proposed quasi-coherent-sheaf replacement. Consequently, these displayed questions retain unresolved qualifications and do not assert a theorem under silently added hypotheses.
\subsection{Short English statement}
The problem asks whether Hochschild cohomology of a heart acts canonically through the graded centre of its ambient triangulated category, whether that action is independent of the heart, and whether an analogous construction exists for schemes and analytic spaces. The meanings of characteristic, comparison between hearts, and several geometric and cohomological conventions remain unspecified. The supplied status is partially solved; this draft does not independently verify that assessment.
\subsection{Sources}
\begin{itemize}
\item Local UnsolvedMath registry: supplied record with permanent ID \(30001189\), exact source problem code \texttt{catalogue entry}, titled \emph{Heart-Independent Hochschild Characteristic Homomorphisms}. The record belongs to the Oberwolfach Reports Open Problems collection and gives proposed year \(2009\); no individual proposer is named.
\item Supplied report citation: \emph{Mini-Workshop: Support Varieties} (2009), Oberwolfach Reports, pp.\ 549--592, DOI \texttt{10.4171/owr/2009/10}. \url{https://doi.org/10.4171/owr/2009/10}. This is also the supplied literature-source link labelled ``original Oberwolfach problem.'' No precise problem page within the report is supplied.
\item Supplied literature source labelled ``Buchweitz--Flenner'': \url{https://arxiv.org/abs/math/0606593}. The record supplies no further bibliographic details or pinpoint citation.
\item Literature assessment dated exactly \texttt{2026-08-21}: the supplied curation labels the problem partially solved and states that the algebra case has a characteristic homomorphism from Hochschild cohomology to the graded centre, that geometric Hochschild theory supplies analogous maps for schemes, and that the general heart-of-a-\(t\)-structure and heart-independence questions remain open. This is supplied dated curation, not an independently verified literature assessment; its time-sensitive status and cited claims require independent verification.
\item The registry labels the statement \texttt{corrected\_verified} and reports checking its mathematical content against the official Oberwolfach report, with notation, context, and question form normalized without changing the source-level claim. This verification assertion is supplied metadata; no source was independently visited or checked for this draft.
\end{itemize}
\end{document}
